\documentclass[11pt,compress,t,notes=noshow, xcolor=table, aspectratio=169]{beamer}
\usepackage[]{graphicx}

\usepackage[]{color}
% maxwidth is the original width if it is less than linewidth
% otherwise use linewidth (to make sure the graphics do not exceed the margin)
\makeatletter
\def\maxwidth{ %
  \ifdim\Gin@nat@width>\linewidth
    \linewidth
  \else
    \Gin@nat@width
  \fi
}
\makeatother

\definecolor{fgcolor}{rgb}{0.345, 0.345, 0.345}
\newcommand{\hlnum}[1]{\textcolor[rgb]{0.686,0.059,0.569}{#1}}%
\newcommand{\hlstr}[1]{\textcolor[rgb]{0.192,0.494,0.8}{#1}}%
\newcommand{\hlcom}[1]{\textcolor[rgb]{0.678,0.584,0.686}{\textit{#1}}}%
\newcommand{\hlopt}[1]{\textcolor[rgb]{0,0,0}{#1}}%
\newcommand{\hlstd}[1]{\textcolor[rgb]{0.345,0.345,0.345}{#1}}%
\newcommand{\hlkwa}[1]{\textcolor[rgb]{0.161,0.373,0.58}{\textbf{#1}}}%
\newcommand{\hlkwb}[1]{\textcolor[rgb]{0.69,0.353,0.396}{#1}}%
\newcommand{\hlkwc}[1]{\textcolor[rgb]{0.333,0.667,0.333}{#1}}%
\newcommand{\hlkwd}[1]{\textcolor[rgb]{0.737,0.353,0.396}{\textbf{#1}}}%
\let\hlipl\hlkwb

\usepackage{framed}
\makeatletter
\newenvironment{kframe}{%
 \def\at@end@of@kframe{}%
 \ifinner\ifhmode%
  \def\at@end@of@kframe{\end{minipage}}%
  \begin{minipage}{\columnwidth}%
 \fi\fi%
 \def\FrameCommand##1{\hskip\@totalleftmargin \hskip-\fboxsep
 \colorbox{shadecolor}{##1}\hskip-\fboxsep
     % There is no \\@totalrightmargin, so:
     \hskip-\linewidth \hskip-\@totalleftmargin \hskip\columnwidth}%
 \MakeFramed {\advance\hsize-\width
   \@totalleftmargin\z@ \linewidth\hsize
   \@setminipage}}%
 {\par\unskip\endMakeFramed%
 \at@end@of@kframe}
\makeatother

\definecolor{shadecolor}{rgb}{.97, .97, .97}
\definecolor{messagecolor}{rgb}{0, 0, 0}
\definecolor{warningcolor}{rgb}{1, 0, 1}
\definecolor{errorcolor}{rgb}{1, 0, 0}
\newenvironment{knitrout}{}{} % an empty environment to be redefined in TeX

\newcommand{\SweaveOpts}[1]{}  % do not interfere with LaTeX
\newcommand{\SweaveInput}[1]{} % because they are not real TeX commands
\newcommand{\Sexpr}[1]{}       % will only be parsed by R
\newcommand{\xmark}{\ding{55}}%

\usepackage{alltt}
\usepackage[utf8]{inputenc}

\usepackage[english]{babel}
\usepackage{dsfont}
\usepackage{verbatim}
\usepackage{amsmath}
\usepackage{amsfonts}
\usepackage{amssymb}
\usepackage{bm}
\usepackage{csquotes}
\usepackage{multirow}
\usepackage{longtable}
\usepackage{booktabs}
\usepackage{enumerate}
\usepackage[absolute,overlay]{textpos}
\usepackage{psfrag}
\usepackage{algorithm}
\usepackage{algpseudocode}
\usepackage{eqnarray}
\usepackage{arydshln}
\usepackage{tabularx}
\usepackage{placeins}
\usepackage{tikz}
\usepackage{setspace}
\usepackage{colortbl}
\usepackage{mathtools}
\usepackage{wrapfig}
\usepackage{bm}
\usepackage{amsmath}
\usepackage{pifont}
\usepackage{xcolor} %colored math symbols
\usepackage{adjustbox}
\usepackage{verbatimbox}
\usepackage{tcolorbox}
\usepackage{siunitx}
\def\signed #1{{\leavevmode\unskip\nobreak\hfil\penalty50\hskip1em
  \hbox{}\nobreak\hfill #1%
  \parfillskip=0pt \finalhyphendemerits=0 \endgraf}}

% basic latex stuff
\newcommand{\col}{\par\colorbox{code}{\parbox{\textwidth}{\theverbbox}}\par}


\usetikzlibrary{shapes,arrows,automata,positioning,calc,chains,trees, shadows}
\tikzset{
  %Define standard arrow tip
  >=stealth',
  %Define style for boxes
  punkt/.style={
    rectangle,
    rounded corners,
    draw=black, very thick,
    text width=6.5em,
    minimum height=2em,
    text centered},
  % Define arrow style
  pil/.style={
    ->,
    thick,
    shorten <=2pt,
    shorten >=2pt,}
}

\newsavebox\mybox
\newenvironment{aquote}[1]
  {\savebox\mybox{#1}\begin{quote}\openautoquote\hspace*{-.7ex}}
  {\unskip\closeautoquote\vspace*{1mm}\signed{\usebox\mybox}\end{quote}}

\usepackage{subfig}

% Defines macros and environments
\usepackage{../style/dhbw-lecture}


\let\code=\texttt
\let\proglang=\textsf

\setkeys{Gin}{width=0.9\textwidth}

% The accent rule sits on its own vbox line with \nointerlineskip so it costs
% only its own 1.2pt height + the explicit 2pt gap, instead of pulling in a
% full \baselineskip (~0.6-0.7cm) the way a plain "\\" line break would --
% that overhead on every single frame is what was pushing picture-heavy
% frames over the shorter 16:9 frame height.
\setbeamertemplate{frametitle}{\expandafter\uppercase\expandafter\insertframetitle\par\vskip2pt\nointerlineskip\hbox{\color{dhbwaccent}\rule{1.6cm}{1.2pt}}}
\addtobeamertemplate{frametitle}{}{\vspace{0.08cm}}

\usepackage{bbm}
% basic latex stuff
\newcommand{\pkg}[1]{{\fontseries{b}\selectfont #1}} %fontstyle for R packages
\newcommand{\lz}{\vspace{0.5cm}} %vertical space
\newcommand{\dlz}{\vspace{1cm}} %double vertical space
\newcommand{\oneliner}[1] % Oneliner for important statements
{\begin{block}{}\begin{center}\begin{Large}#1\end{Large}\end{center}\end{block}}


%new environments
\newenvironment{vbframe}  %frame with breaks and verbatim
{
 \begin{frame}[containsverbatim,allowframebreaks]
}
{
\end{frame}
}

\newenvironment{vframe}  %frame with verbatim without breaks (to avoid numbering one slided frames)
{
 \begin{frame}[containsverbatim]
}
{
\end{frame}
}

\newenvironment{blocki}[1]   % itemize block
{
 \begin{block}{#1}\begin{itemize}
}
{
\end{itemize}\end{block}
}

\newenvironment{fragileframe}[2]{  %fragile frame with framebreaks
\begin{frame}[allowframebreaks, fragile, environment = fragileframe]
\frametitle{#1}
#2}
{\end{frame}}


\newcommand{\myframe}[2]{  %short for frame with framebreaks
\begin{frame}[allowframebreaks]
\frametitle{#1}
#2
\end{frame}}

\newcommand{\remark}[1]{
  \textbf{Remark:} #1
}


\newenvironment{deleteframe}
{
\begingroup
\usebackgroundtemplate{\includegraphics[width=\paperwidth,height=\paperheight]{../style/color/red.png}}
 \begin{frame}
}
{
\end{frame}
\endgroup
}
\newenvironment{simplifyframe}
{
\begingroup
\usebackgroundtemplate{\includegraphics[width=\paperwidth,height=\paperheight]{../style/color/yellow.png}}
 \begin{frame}
}
{
\end{frame}
\endgroup
}\newenvironment{draftframe}
{
\begingroup
\usebackgroundtemplate{\includegraphics[width=\paperwidth,height=\paperheight]{../style/color/green.jpg}}
 \begin{frame}
}
{
\end{frame}
\endgroup
}
% https://tex.stackexchange.com/a/261480: textcolor that works in mathmode
\makeatletter
\renewcommand*{\@textcolor}[3]{%
  \protect\leavevmode
  \begingroup
    \color#1{#2}#3%
  \endgroup
}
\makeatother


\input{../latex-math/basic-math}
\input{../latex-math/basic-ml}
\input{../latex-math/ml-nn}

% The course's typewriter font (laett) has no bold shape; map bold typewriter to regular.
\AtBeginDocument{{\ttfamily\selectfont}\DeclareFontShape{T1}{laett}{bx}{n}{<->ssub*laett/m/n}{}\DeclareFontShape{T1}{laett}{m}{bx}{<->ssub*laett/m/n}{}}
\newcommand{\citebutton}[2]{\beamergotobutton{\href{#2}{#1}}}

\newcommand{\titlefigure}{figure/title-modern-variant.png}
\newcommand{\learninggoals}{
  \item What changed in the Transformer block since 2017 and why: pre-norm, RMSNorm, no biases, gated activations
  \item Rotary position embeddings (RoPE)
  \item The consensus hyperparameters of modern language models
  \item Stability tricks and attention variants (MQA/GQA, sliding windows), and why memory matters
  \item Mixture of experts: routing, training, DeepSeek-MoE
}

\title{Advanced Artificial Intelligence and Machine Learning}
\date{}

\begin{document}

\lecturechapter{Transformers Modernized}
\lecture{Advanced AI \& ML}

%%%%%%%%%%%%%%%%%%%%%%%%%%%%%%%%%%%%%%%%%%%%%%%%%%%%%%%%%%%%%%%%%%
\section{Since 2017: What Changed Inside the Transformer?}
% Source: CS336 Spring 2025, lecture 3 (architecture), slides 3-8
%%%%%%%%%%%%%%%%%%%%%%%%%%%%%%%%%%%%%%%%%%%%%%%%%%%%%%%%%%%%%%%%%%

% [O] original bridge frame (draft for the user to rewrite); listed in order.txt
\begin{frame}{Since 2017: What Changed Inside the Transformer?}
\begin{itemize}
  \item Session 7 introduced the original 2017 block: post-norm with LayerNorm, sinusoidal position embeddings, a ReLU feed-forward network. GPT, BERT and T5 (session 8) already changed some of these details, e.g.\ learned or relative position embeddings.
  \item Today's large models (e.g.\ the LLaMA family, DeepSeek) keep the overall design, but change several components \emph{inside} the block.
  \item This session goes through these changes and the evidence behind them: normalization, activations, position embeddings, hyperparameters, stability, attention heads.
  \item Then a change at the level of the whole layer: \textbf{mixture of experts}.
  \item Source: Stanford CS336, Lectures 3 and 4 (Spring 2025).
\end{itemize}
\end{frame}

% FRAMES-BEGIN l3 3 8
% CS336 slide 3
\begin{frame}{Outline and goals}
\begin{itemize}
  \item Quick recap of the `standard' transformer
  \begin{itemize}
    \item What do most of the large LMs have in common?
  \end{itemize}
  \item What are common variations to the architecture / training process?
  \begin{itemize}
    \item Today's theme: the best way to learn is hands-on experience
    \item the second best way is to try to learn from others' experience
  \end{itemize}
\end{itemize}
\end{frame}

% CS336 slide 4
\begin{frame}{Starting point: the `original' transformer}
% hand-typeset (_generator/handwritten/)
\begin{columns}[T]
\begin{column}{0.4\textwidth}
\centering
\includegraphics[page=4,trim=79 5 385 68,clip,width=\linewidth,height=165pt,keepaspectratio]{figure/cs336-l3.pdf}
\end{column}
\begin{column}{0.58\textwidth}
\small
\textbf{Review:} choices in the standard transformer

\medskip
\textbf{Position embedding:} sines and cosines
\begin{align*}
PE_{(pos,2i)} &= \sin\left(pos/10000^{2i/d_{\text{model}}}\right)\\
PE_{(pos,2i+1)} &= \cos\left(pos/10000^{2i/d_{\text{model}}}\right)
\end{align*}
\textbf{FFN:} ReLU
\[ \text{FFN}(x) = \max(0, xW_1 + b_1)W_2 + b_2 \]
\textbf{Norm type:} post-norm, LayerNorm
\end{column}
\end{columns}
\end{frame}

% CS336 slide 5
\begin{frame}{A Simple, Modern Variant}
% hand-typeset (_generator/handwritten/)
\begin{columns}[T]
\begin{column}{0.45\textwidth}
\centering
\includegraphics[page=5,trim=24 58 363 100,clip,width=\linewidth,height=150pt,keepaspectratio]{figure/cs336-l3.pdf}
\end{column}
\begin{column}{0.53\textwidth}
\textbf{Differences:}
\begin{itemize}
  \item \textbf{LayerNorm} is in front of the block
  \item \textbf{Rotary position embeddings (RoPE)}
  \item FF layers use \textbf{SwiGLU}, not ReLU
  \item Linear layers (and layernorm) have \textbf{no bias} (constant) terms
\end{itemize}
\bigskip
\centering
\textbf{Why did we pick these?}\\
\textbf{What should you pick?}
\end{column}
\end{columns}
\end{frame}

% CS336 slide 6
\begin{frame}{How should we think about architectures?}
% hand-typeset (_generator/handwritten/)
Lots of architecture. Just in the last year since last 336..
\begin{center}
\includegraphics[page=6,trim=31 195 47 95,clip,width=250pt]{figure/cs336-l3.pdf}\\[2pt]
\colorbox{blue!8}{\small Over 19 new \emph{dense} model releases, many of them with minor architecture tweaks..}\\[2pt]
\includegraphics[page=6,trim=31 11 47 234,clip,width=250pt]{figure/cs336-l3.pdf}
\end{center}
\end{frame}

% CS336 slide 7
\begin{frame}{Let's look at the data (on dense architectures)}
% hand-typeset (_generator/handwritten/)
Learn from the many other models (and papers) out there
\begin{columns}[T]
\begin{column}{0.5\textwidth}
\centering
\includegraphics[page=7,trim=72 8 328 101,clip,width=\linewidth,height=145pt,keepaspectratio]{figure/cs336-l3.pdf}
\end{column}
\begin{column}{0.48\textwidth}
\vspace{2.5em}
\textbf{We will talk through many major architecture and hyperparameter variants.}
\begin{itemize}
  \item What do all these models have in common?
  \item What parts vary?
  \item What can we learn from this?
\end{itemize}
\end{column}
\end{columns}
\end{frame}

% CS336 slide 8
\begin{frame}{What are we going to cover?}
\begin{itemize}
  \item Common architecture variations
  \begin{itemize}
    \item Activations, FFN
    \item Attention variants
    \item Position embeddings
  \end{itemize}
  \item Hyperparameters that (do or don't) matter
  \begin{itemize}
    \item What is ff\_dim? Do multi\_head dims always sum to model\_dim?
    \item How many vocab elements?
  \end{itemize}
  \item Stability tricks
\end{itemize}
\end{frame}

% FRAMES-END

%%%%%%%%%%%%%%%%%%%%%%%%%%%%%%%%%%%%%%%%%%%%%%%%%%%%%%%%%%%%%%%%%%
\section{Architecture Variations: Normalization}
% Source: CS336 Spring 2025, lecture 3 (architecture), slides 9-14
%%%%%%%%%%%%%%%%%%%%%%%%%%%%%%%%%%%%%%%%%%%%%%%%%%%%%%%%%%%%%%%%%%

% FRAMES-BEGIN l3 9 14
% CS336 slide 9
\begin{frame}{Architecture variations..}
% hand-typeset (_generator/handwritten/)
\begin{itemize}
  \item Let's think about the core architecture piece
  \begin{itemize}
    \item Trends toward `LLaMA- like' architectures
  \end{itemize}
\end{itemize}
\begin{columns}[T]
\begin{column}{0.64\textwidth}
\centering
\includegraphics[page=9,trim=20 10 259 98,clip,width=\linewidth,height=125pt,keepaspectratio]{figure/cs336-l3.pdf}
\end{column}
\begin{column}{0.34\textwidth}
\vspace{1.5em}
\textbf{High level view:}
\begin{itemize}
  \item Low consensus (except pre-norm)
  \item Trends toward `LLaMA-like' architectures
\end{itemize}
\end{column}
\end{columns}
\end{frame}

% CS336 slide 10
\begin{frame}{Pre-vs-post norm}
% hand-typeset (_generator/handwritten/)
The one thing \emph{everyone} agrees on (in 2024)
\begin{columns}[T]
\begin{column}{0.36\textwidth}
\centering
\includegraphics[page=10,trim=51 75 441 103,clip,width=\linewidth,height=95pt,keepaspectratio]{figure/cs336-l3.pdf}\\
{\scriptsize Figure from Xiong 2020\par}
\end{column}
\begin{column}{0.62\textwidth}
\centering
\includegraphics[page=10,trim=290 154 20 127,clip,width=\linewidth]{figure/cs336-l3.pdf}\par
\medskip
\small Set up LayerNorm so that it doesn't affect the main residual signal path (on the left)
\end{column}
\end{columns}
\medskip
{\small\centering
\textbf{Almost all modern LMs use pre-norm (but BERT was post-norm)}\\
(One somewhat funny exception -- OPT350M. I don't know why this is post-norm)\par}
\end{frame}

% CS336 slide 11
\begin{frame}{Pre-vs-post-norm, the data}
% hand-typeset (_generator/handwritten/)
\begin{columns}[T]
\begin{column}{0.48\textwidth}
\centering
\includegraphics[page=11,trim=29 58 378 93,clip,width=\linewidth,height=150pt,keepaspectratio]{figure/cs336-l3.pdf}\\
{\scriptsize Salazar and Ngyuen 2019\par}
\end{column}
\begin{column}{0.5\textwidth}
\centering
\includegraphics[page=11,trim=363 178 -1 68,clip,width=\linewidth]{figure/cs336-l3.pdf}\\
\includegraphics[page=11,trim=374 18 190 221,clip,width=\linewidth,height=72pt,keepaspectratio]{figure/cs336-l3.pdf}\\
{\scriptsize Figure from Xiong 2020\par}
\end{column}
\end{columns}
\end{frame}

% CS336 slide 12
\begin{frame}{Pre-vs-post norm, explanations?}
% hand-typeset (_generator/handwritten/)
\begin{columns}[T]
\begin{column}{0.49\textwidth}
\centering
Gradient attenuation [Xiong 2020]\\[2pt]
\includegraphics[page=12,trim=85 80 369 104,clip,width=\linewidth,height=105pt,keepaspectratio]{figure/cs336-l3.pdf}
\end{column}
\begin{column}{0.49\textwidth}
\centering
Gradient spikes [Salazar and Ngyuen]\\[2pt]
\includegraphics[page=12,trim=392 82 44 104,clip,width=\linewidth,height=105pt,keepaspectratio]{figure/cs336-l3.pdf}
\end{column}
\end{columns}
\medskip
{\centering
\textbf{Original stated advantage} -- removing warmup.\\
\textbf{Today} -- stability and larger LRs for large networks\par}
\end{frame}

% CS336 slide 13
\begin{frame}{New things -- `double' norm.}
\begin{itemize}
  \item If putting LayerNorms in residual streams is bad.. Why not post-norm outside the stream?
  \item Recent models: Grok, Gemma 2. Olmo 2 only does non-residual post norm
\end{itemize}
\begin{center}
  \includegraphics[page=13,trim=246 56 246 118,clip,width=119pt]{figure/cs336-l3.pdf}\par
\end{center}
\end{frame}

% CS336 slide 14
\begin{frame}{LayerNorm vs RMSNorm}
% hand-typeset (_generator/handwritten/)
\small
\renewcommand{\arraystretch}{1.3}
\makebox[\textwidth][c]{%
\begin{tabular}{@{}>{\raggedright\arraybackslash}p{150pt}@{\hspace{12pt}}>{\centering\arraybackslash}m{130pt}@{\hspace{12pt}}>{\raggedright\arraybackslash}p{95pt}@{}}
\toprule
 & & \textbf{Notable models}\\
\midrule
Original transformer: \textbf{LayerNorm} -- normalizes the mean and variance across $d_{\text{model}}$
 & $\displaystyle y = \frac{x - \operatorname{E}[x]}{\sqrt{\operatorname{Var}[x] + \epsilon}} * \gamma + \beta$
 & GPT3/2/1, OPT, GPT-J, BLOOM\\
\midrule
Many modern LMs: \textbf{RMSNorm} -- does not subtract mean or add a bias term
 & $\displaystyle y = \frac{x}{\sqrt{\|x\|_2^2 + \varepsilon}} * \gamma$
 & LLaMA-family, PaLM, Chinchilla, T5\\
\bottomrule
\end{tabular}}
\end{frame}

% FRAMES-END

% [O] original bridge frame (draft for the user to rewrite); listed in order.txt
\begin{frame}{Compute vs.\ Memory}
\begin{itemize}
  \item A GPU can do arithmetic much faster than it can move data between its memory and its compute units.
  \item The number of arithmetic operations per byte moved is an operation's \textbf{arithmetic intensity}. With low arithmetic intensity, an operation is limited by \textbf{memory traffic}, not by the number of floating-point operations (FLOPs).
  \item Matrix multiplications do many operations per byte loaded; normalizations and element-wise operations do few.
  \item So an operation can account for a tiny share of the FLOPs and still a noticeable share of the runtime --- the next slides use exactly this argument.
  \item Hardware details follow in session 12.
\end{itemize}
\end{frame}

% FRAMES-BEGIN l3 15 19
% CS336 slide 15
\begin{frame}{Why RMSNorm?}
\begin{itemize}
  \item Modern explanation -- it's faster (and just as good).
  \begin{itemize}
    \item Fewer operations (no mean calculation)
    \item Fewer parameters (no bias term to store)
  \end{itemize}
  \item Does this explanation make sense?
  \begin{itemize}
    \item Matrix multiplies are the vast majority of FLOPs (and memory)
  \end{itemize}
\end{itemize}
\begin{center}
  \includegraphics[page=15,trim=464 214 23 112,clip,width=95pt]{figure/cs336-l3.pdf}\par
  \includegraphics[page=15,trim=247 61 250 244,clip,width=91pt]{figure/cs336-l3.pdf}\par
\end{center}
{\scriptsize [Ivanov et al 2023]\par}
\end{frame}

% CS336 slide 16
\begin{frame}[shrink=7]{Why RMSNorm (2)}
\begin{itemize}
  \item Important lesson: FLOPS are not runtime!
  \begin{itemize}
    \item RMSNorm can still matter due to the importance of data movement
  \end{itemize}
\end{itemize}
\begin{center}
  \includegraphics[page=16,trim=470 79 143 107,clip,width=42pt]{figure/cs336-l3.pdf}\par
  \includegraphics[page=16,trim=100 185 336 130,clip,width=111pt]{figure/cs336-l3.pdf}\par
\end{center}
{\scriptsize Left top (''43G'') is FLOPS Right top (``153'') is the FLOP-to-memory ratio; [Ivanov et al 2023]\par}
\end{frame}

% CS336 slide 17
\begin{frame}{RMSNorm - validation}
% hand-typeset (_generator/handwritten/)
\begin{itemize}
  \item \textbf{RMSNorm} runtime (and surprisingly, perf) gains have been seen in papers
\end{itemize}
\vspace{1.5em}
\begin{center}
\includegraphics[page=17,trim=59 124 43 169,clip,width=\linewidth]{figure/cs336-l3.pdf}
\end{center}
{\raggedleft\scriptsize Narang et al 2020\par}
\end{frame}

% CS336 slide 18
\begin{frame}{More generally: dropping bias terms}
% hand-typeset (_generator/handwritten/)
\textbf{Most modern transformers don't have bias terms.}

\bigskip
Original Transformer:
\[ \text{FFN}(x) = \max(0, xW_1 + b_1)\,W_2 + b_2 \]
Most implementations (if they're not gated):
\[ FFN(x) = \sigma(xW_1)\,W_2 \]
\medskip
{\centering \textbf{Reasons:} memory (similar to RMSnorm) and optimization stability\par}
\end{frame}

% CS336 slide 19
\begin{frame}{LayerNorm: recap}
\begin{itemize}
  \item Basically everyone does pre-norm.
  \item Intuition -- keep the good parts of residual connections
  \item Observations -- nicer gradient propagation, fewer spike
  \item Some people add a second norm outside the residual stream (NOT post-norm)
  \item Most people do RMSnorm
  \item In practice, works as well as LayerNorm
  \item But, has fewer parameters to move around, which saves on wallclock time
  \item People more generally drop bias terms since the compute/param tradeoffs are not great.
\end{itemize}
\end{frame}

% FRAMES-END

%%%%%%%%%%%%%%%%%%%%%%%%%%%%%%%%%%%%%%%%%%%%%%%%%%%%%%%%%%%%%%%%%%
\section{Architecture Variations: Activations and Layers}
% Source: CS336 Spring 2025, lecture 3 (architecture), slides 20-29
%%%%%%%%%%%%%%%%%%%%%%%%%%%%%%%%%%%%%%%%%%%%%%%%%%%%%%%%%%%%%%%%%%

% FRAMES-BEGIN l3 20 29
% CS336 slide 20
\begin{frame}{Activations}
\begin{itemize}
  \item A whole zoo of activations ..
  \item ReLU, GeLU, Swish, ELU, GLU, GeGLU, ReGLU, SeLU, SwiGLU, LiGLU
  \begin{itemize}
    \item What are these things? What do people use? Does it matter?
  \end{itemize}
\end{itemize}
\end{frame}

% CS336 slide 21
\begin{frame}{A few of the common activations}
% hand-typeset (_generator/handwritten/)
\small
\renewcommand{\arraystretch}{1.2}
\makebox[\textwidth][c]{%
\begin{tabular}{@{}>{\raggedright\arraybackslash}m{40pt}@{\hspace{8pt}}>{\raggedright\arraybackslash}m{120pt}@{\hspace{8pt}}>{\centering\arraybackslash}m{90pt}@{\hspace{12pt}}>{\raggedright\arraybackslash}m{110pt}@{}}
\toprule
 & & & \textbf{Notable models}\\
\midrule
\textbf{ReLU} & $FF(x) = \max(0, xW_1)\,W_2$
 & \includegraphics[page=21,trim=357 231 218 64,clip,width=\linewidth,height=52pt,keepaspectratio]{figure/cs336-l3.pdf}
 & Original transformer, T5, Gopher, Chinchilla, OPT\\
\midrule
\textbf{GeLU} & $FF(x) = \text{GELU}(xW_1)\,W_2$\newline $\text{GELU}(x) := x\,\Phi(x)$
 & \includegraphics[page=21,trim=356 107 219 187,clip,width=\linewidth,height=52pt,keepaspectratio]{figure/cs336-l3.pdf}
 & GPT1/2/3, GPTJ, GPT-Neox, BLOOM\\
\midrule
\multicolumn{3}{@{}l}{\textbf{SwiGLU / GeGLU} (next slide..)}
 & Llama, PaLM, T5 v1.1, \emph{most models post 2023}\\
\bottomrule
\end{tabular}}
\end{frame}

% CS336 slide 22
\begin{frame}{Gated activations (*GLU)}
% hand-typeset (_generator/handwritten/)
GLUs modify the `first part' of a FF layer
\[ FF(x) = \textcolor{red}{\max(0, xW_1)}\,W_2 \]
Instead of a linear + ReLU, augment the above with an (entrywise) linear term
\[ \textcolor{red}{\max(0, xW_1)} \rightarrow \max(0, xW_1) \otimes (xV) \]
This gives the gated variant (ReGLU) -- note that we have an extra parameter ($V$)
\[ \text{FF}_{\text{ReGLU}}(x) = (\max(0, xW_1) \otimes xV)\,W_2 \]
\end{frame}

% CS336 slide 23
\begin{frame}{Gated variants of standard FF layers}
% hand-typeset (_generator/handwritten/)
\small
\renewcommand{\arraystretch}{1.3}
\makebox[\textwidth][c]{%
\begin{tabular}{@{}>{\raggedright\arraybackslash}p{270pt}@{\hspace{12pt}}>{\raggedright\arraybackslash}p{110pt}@{}}
\toprule
 & \textbf{Notable models}\\
\midrule
\textbf{GeGLU}\newline
$\text{FFN}_{\text{GEGLU}}(x, W, V, W_2) = (\text{GELU}(xW) \otimes xV)\,W_2$
 & T5 v1.1, mT5, LaMDA, Phi3, Gemma 2, Gemma 3\\
\midrule
\textbf{SwiGLU} (swish is $x * \text{sigmoid}(x)$)\newline
$\text{FFN}_{\text{SwiGLU}}(x, W, V, W_2) = (\text{Swish}_1(xW) \otimes xV)\,W_2$
 & LLaMa 1/2/3, PaLM, Mistral, OlMo, \emph{most models post 2023}\\
\bottomrule
\end{tabular}}

\bigskip
{\centering Note: Gated models use smaller dimensions for the $d_{ff}$ by 2/3\par}
\end{frame}

% CS336 slide 24
\begin{frame}{Do gated linear units work?}
% hand-typeset (_generator/handwritten/)
Yes, fairly consistently so.
\begin{center}
\includegraphics[page=24,trim=225 75 188 109,clip,width=\linewidth,height=135pt,keepaspectratio]{figure/cs336-l3.pdf}
\end{center}
{\raggedleft\scriptsize Shazeer 2020\par}
\end{frame}

% CS336 slide 25
\begin{frame}{Do gated linear units work (2)?}
\begin{itemize}
  \item Yes, with other works corroborating Shazeer 2020
  \begin{itemize}
    \item Narang et al 2020
  \end{itemize}
\end{itemize}
\begin{center}
  \includegraphics[page=25,trim=80 95 51 105,clip,width=330pt]{figure/cs336-l3.pdf}\par
\end{center}
\end{frame}

% CS336 slide 26
\begin{frame}{Gating, activations}
\begin{itemize}
  \item Many variations (ReLU, GeLU, *GLU) across models.
  \item *GLU isn't necessary for a good model (see GPT3), but it's probably helpful Otuer, recent outlier models.. Nemotron 340B (Squared ReLU), Falcon 2 11b (ReLU)
  \item But evidence points towards somewhat consistent gains from Swi/GeGLU
\end{itemize}
\end{frame}

% CS336 slide 27
\begin{frame}{Serial vs Parallel layers}
\begin{itemize}
  \item Normal transformer blocks are serial -- they compute attention, then the MLP
  \begin{itemize}
    \item Could we parallelize the transformer block?
  \end{itemize}
\end{itemize}
\begin{center}
  \includegraphics[page=27,trim=284 83 284 103,clip,width=85pt]{figure/cs336-l3.pdf}\par
\end{center}
\end{frame}

% CS336 slide 28
\begin{frame}{Parallel layers}
\begin{itemize}
  \item A few models (GPTJ, PaLM, GPT-NeoX) do parallel layers. Originally in GPT-J
  \item If implemented right, LayerNorm can be shared, and matrix multiplies can be fused
  \item Recent Models: Cohere Command A, Falcon 2 11B, Command R+
\end{itemize}
\begin{center}
  \includegraphics[page=28,trim=91 96 55 115,clip,width=310pt]{figure/cs336-l3.pdf}\par
\end{center}
\end{frame}

% CS336 slide 29
\begin{frame}{Summary: architectures}
% hand-typeset (_generator/handwritten/)
\makebox[\textwidth][c]{%
\begin{minipage}{400pt}
\begin{columns}[T,onlytextwidth]
\begin{column}{0.58\textwidth}
\footnotesize
\setlength{\itemsep}{0pt}
\textbf{Pre-vs-post norm:}
\begin{itemize}
  \item Everyone does pre-norm (except OPT350M), likely with good reason.
\end{itemize}
\textbf{Layer vs RMSnorm:}
\begin{itemize}
  \item RMSnorm has clear compute wins, sometimes even performance
\end{itemize}
\textbf{Gating:}
\begin{itemize}
  \item GLUs seem generally better, though differences are small
\end{itemize}
\textbf{Serial vs parallel layers:}
\begin{itemize}
  \item No extremely serious ablations, but has a compute win.
\end{itemize}
\end{column}
\begin{column}{0.4\textwidth}
\centering
\includegraphics[page=29,trim=406 20 42 35,clip,width=\linewidth,height=170pt,keepaspectratio]{figure/cs336-l3.pdf}
\end{column}
\end{columns}
\end{minipage}}
\end{frame}

% FRAMES-END

%%%%%%%%%%%%%%%%%%%%%%%%%%%%%%%%%%%%%%%%%%%%%%%%%%%%%%%%%%%%%%%%%%
\section{Position Embeddings: RoPE}
% Source: CS336 Spring 2025, lecture 3 (architecture), slides 30-35
%%%%%%%%%%%%%%%%%%%%%%%%%%%%%%%%%%%%%%%%%%%%%%%%%%%%%%%%%%%%%%%%%%

% FRAMES-BEGIN l3 30 35
% CS336 slide 30
\begin{frame}{Many variations in position embeddings}
% hand-typeset (_generator/handwritten/)
\footnotesize
\renewcommand{\arraystretch}{1.2}
\makebox[\textwidth][c]{%
\begin{tabular}{@{}>{\raggedright\arraybackslash}p{130pt}@{\hspace{10pt}}>{\centering\arraybackslash}p{160pt}@{\hspace{10pt}}>{\raggedright\arraybackslash}p{100pt}@{}}
\toprule
 & & \textbf{Notable models}\\
\midrule
\textbf{Sine embeddings:} add sines and cosines that enable localization
 & $\mathit{Embed}(x, i) = v_x + PE_{pos}$\newline
   $PE_{(pos,2i)} = \sin\left(pos/10000^{2i/d_{\text{model}}}\right)$\newline
   $PE_{(pos,2i+1)} = \cos\left(pos/10000^{2i/d_{\text{model}}}\right)$
 & Original transformer\\
\midrule
\textbf{Absolute embeddings:} add a position vector to the embedding
 & $\mathit{Embed}(x, i) = v_x + u_i$
 & GPT1/2/3, OPT\\
\midrule
\textbf{Relative embeddings:} add a vector to the \emph{attention computation}
 & $e_{ij} = \dfrac{x_i W^Q \left(x_j W^K + a_{ij}^K\right)^T}{\sqrt{d_z}}$
 & T5, Gopher, Chinchilla\\
\midrule
\textbf{Rope embeddings} (next slides..)
 &
 & GPTJ, PaLM, LLaMA,\newline \emph{Most 2024+ models}\\
\bottomrule
\end{tabular}}
\end{frame}

% CS336 slide 31
\begin{frame}{RoPE: rotary position embeddings}
% hand-typeset (_generator/handwritten/)
\small
\textbf{High level thought process:} a \emph{relative} position embedding should be some $f(x, i)$ s.t.
\[ \langle f(x, i), f(y, j) \rangle = g(x, y, i - j) \]
That is, the attention function \emph{only} gets to depend on the relative position ($i-j$). How do existing embeddings not fulfill this goal?
\begin{itemize}
  \item \textbf{Sine:} Has various cross-terms that are not relative
  \[ \langle \mathit{Embed}(x, i), \mathit{Embed}(y, i) \rangle = \langle v_x, v_y \rangle + \langle PE_i, v_y \rangle \ldots \]
  \item \textbf{Absolute:} obviously not relative
  \item \textbf{Relative embeddings:}
  \[ e_{ij} = \frac{x_i W^Q \left(x_j W^K + a_{ij}^K\right)^T}{\sqrt{d_z}} \quad \text{is not an inner product} \]
\end{itemize}
\end{frame}

% CS336 slide 32
\begin{frame}{RoPE: rotary position embeddings}
% hand-typeset (_generator/handwritten/)
\textbf{How can we solve this problem?}
\begin{itemize}
  \item We want our embeddings to be \textbf{invariant to absolute position}
  \item We know that inner products are invariant to arbitrary rotation.
\end{itemize}
\medskip
\footnotesize
\renewcommand{\arraystretch}{1.1}
\makebox[\textwidth][c]{%
\begin{tabular}{@{}*{3}{>{\centering\arraybackslash}p{125pt}}@{}}
\includegraphics[page=32,trim=122 107 472 169,clip,width=\linewidth,height=58pt,keepaspectratio]{figure/cs336-l3.pdf} & \includegraphics[page=32,trim=342 105 280 169,clip,width=\linewidth,height=58pt,keepaspectratio]{figure/cs336-l3.pdf} & \includegraphics[page=32,trim=467 112 89 192,clip,width=\linewidth,height=45pt,keepaspectratio]{figure/cs336-l3.pdf}\\[4pt]
\textbf{Position independent embedding} & \textbf{Embedding ``we know that''} & \textbf{Embedding ``of course we know''}\\
 & Rotate we by `0 positions'\newline know by `1 positions' & Rotate we by `2 positions'\newline Rotate know by `3 positions'
\end{tabular}}
\end{frame}

% CS336 slide 33
\begin{frame}{RoPE: rotary position embeddings}
\begin{itemize}
  \item There are many rotations, which one do you pick?
  \begin{itemize}
    \item Just pair up the coordinates and rotate them in 2d (motivation: complex numbers)
  \end{itemize}
\end{itemize}
\begin{center}
  \includegraphics[page=33,trim=131 53 148 99,clip,width=211pt]{figure/cs336-l3.pdf}\par
\end{center}
{\scriptsize [Su et al 2021]\par}
\end{frame}

% CS336 slide 34
\begin{frame}{The actual RoPE math}
\begin{itemize}
  \item Multiply with sines and cosines
  \item Difference with sine embeddings -- not additive, no cross terms
\end{itemize}
\begin{center}
  \includegraphics[page=34,trim=80 111 72 122,clip,width=330pt]{figure/cs336-l3.pdf}\par
\end{center}
\end{frame}

% CS336 slide 35
\begin{frame}{Implementation and code for RoPE}
% hand-typeset (_generator/handwritten/)
\makebox[\textwidth][c]{\begin{minipage}{400pt}
{\footnotesize\textcolor{dhbwink}{$\triangleright$ Usual attention stuff}\par}
{\ttfamily\fontsize{6.4}{7.6}\selectfont query\_states = self.q\_proj(hidden\_states)\\
key\_states = self.k\_proj(hidden\_states)\\
value\_states = self.v\_proj(hidden\_states)\\
\mbox{}\\
\# Flash attention requires the input to have the shape\\
\# batch\_size x seq\_length x head\_dim x hidden\_dim\\
\# therefore we just need to keep the original shape\\
query\_states = query\_states.view(bsz, q\_len, self.num\_heads, self.head\_dim).transpose(1, 2)\\
key\_states = key\_states.view(bsz, q\_len, self.num\_key\_value\_heads, self.head\_dim).transpose(1, 2)\\
value\_states = value\_states.view(bsz, q\_len, self.num\_key\_value\_heads, self.head\_dim).transpose(1, 2)\par}
\smallskip
{\footnotesize\textcolor{dhbwink}{$\triangleright$ Get the RoPE matrix cos/sin}\par}
{\ttfamily\fontsize{6.4}{7.6}\selectfont cos, sin = self.rotary\_emb(value\_states, position\_ids)\par}
\smallskip
{\footnotesize\textcolor{dhbwink}{$\triangleright$ Multiply query/key inputs}\par}
{\ttfamily\fontsize{6.4}{7.6}\selectfont query\_states, key\_states = apply\_rotary\_pos\_emb(query\_states, key\_states, cos, sin)\par}
{\centering\footnotesize \ldots\ Same stuff as the usual multi-head self attention below\par}
\end{minipage}}\par
\vspace{8pt}
{\centering \textbf{Note:} embedding at \emph{each attention operation} to enforce position invariance\par}
\end{frame}

% FRAMES-END

%%%%%%%%%%%%%%%%%%%%%%%%%%%%%%%%%%%%%%%%%%%%%%%%%%%%%%%%%%%%%%%%%%
\section{Hyperparameters}
% Source: CS336 Spring 2025, lecture 3 (architecture), slides 36-52
%%%%%%%%%%%%%%%%%%%%%%%%%%%%%%%%%%%%%%%%%%%%%%%%%%%%%%%%%%%%%%%%%%

% FRAMES-BEGIN l3 36 52
% CS336 slide 36
\begin{frame}{Hyperparameters}
\begin{itemize}
  \item Transformer hyperparameter questions you might have had
  \item How much bigger should the feedforward size be compared to hidden size?
  \item How many heads, and should num\_heads always divide hidden size?
  \item What should my vocab size be?
  \item And other model setting questions
  \begin{itemize}
    \item Do people even regularize these huge LMs?
    \item How do people scale these models - very deep or very wide?
  \end{itemize}
\end{itemize}
\end{frame}

% CS336 slide 37
\begin{frame}{Surprising (?) consensus hyperparameter 1}
% hand-typeset (_generator/handwritten/)
Feedforward -- model dimension ratio.
\[ \text{FFN}(x) = \max(0, xW_1 + b_1)\,W_2 + b_2 \]
There are two dimensions that are relevant -- the feedforward dim ($d_{ff}$) and model dim ($d_{model}$). What should their relationship be?
\[ \boldsymbol{d_{ff} = 4\, d_{model}} \]
This is \emph{almost always} true. There's just a few exceptions.
\end{frame}

% CS336 slide 38
\begin{frame}{Exception \#1 -- GLU variants}
% hand-typeset (_generator/handwritten/)
\small
Remember that GLU variants scale down by 2/3rd. This means most GLU variants have $d_{ff} = \frac{8}{3} d_{model}$. This is mostly what happens. Some notable such examples.

\begin{center}
\footnotesize
\renewcommand{\arraystretch}{0.9}
\begin{tabular}{@{}lr@{}}
\toprule
\textbf{Model} & $\boldsymbol{d_{ff}/d_{model}}$\\
\midrule
PaLM & 4\\
Mistral 7B & 3.5\\
LLaMA-2 70B & 3.5\\
LLaMA 70B & 2.68\\
Qwen 14B & 2.67\\
DeepSeek 67B & 2.68\\
Yi 34B & 2.85\\
T5 v1.1 & 2.5\\
\bottomrule
\end{tabular}
\end{center}
Models are roughly in this range, though PaLM, LLaMA2 and Mistral are slightly larger
\end{frame}

% CS336 slide 39
\begin{frame}{Exception \#2 -- T5}
% hand-typeset (_generator/handwritten/)
\small
As we have (and will) see, most LMs are have boring, conservative hyperparameters.
One exception is T5 [Raffel et al 2020] which has some \emph{very bold} settings.

\smallskip
In particular, for the 11B model, they set
\[ d_{ff} = \text{65,536} \qquad d_{model} = 1024 \]
For an astounding 64-times multiplier.
\begin{center}
\includegraphics[page=39,trim=116 34 116 276,clip,width=0.8\linewidth]{figure/cs336-l3.pdf}
\end{center}
Other, recent exceptions -- Gemma 2 (8x), SmolLM/Gemma 3 (4x, GLU)
\end{frame}

% CS336 slide 40
\begin{frame}{Why this range of multipliers?}
\begin{itemize}
  \item Empirically, there's a basin between 1-10 where this hyperparameter is near-optimal
\end{itemize}
\begin{center}
  \includegraphics[page=40,trim=184 62 212 104,clip,width=185pt]{figure/cs336-l3.pdf}\par
\end{center}
{\scriptsize [Kaplan+ 2020]\par}
\end{frame}

% CS336 slide 41
\begin{frame}{What can we learn from the model-dim hyperparam?}
% hand-typeset (_generator/handwritten/)
\begin{itemize}
  \item The `default' choices of $d_{ff} = 4 d_{model}$ and $d_{ff} = 2.66 d_{model}$ have worked well for nearly all modern LLMs.
  \medskip
  \item But T5 does show that even radical choices of $d_{ff} = 64 d_{model}$ can work. This hyperparameter choice isn't written in stone.
  \medskip
  \item That said, T5 has a follow-up model (T5 v1.1) that is `improved' and uses a much more standard 2.5 multiplier on GeGLU, so the 64-times multiplier is likely suboptimal.
\end{itemize}
\end{frame}

% CS336 slide 42
\begin{frame}{Surprising (?) consensus hyperparameter 2}
\begin{itemize}
  \item Head-dim*num-heads to model-dim ratio.
  \item This doesn't have to be true: we can have head-dimensions > model-dim / num-heads.
  \begin{itemize}
    \item But most models do follow this guideline
  \end{itemize}
\end{itemize}
\begin{center}
  \includegraphics[page=42,trim=78 142 78 116,clip,width=330pt]{figure/cs336-l3.pdf}\par
\end{center}
\end{frame}

% CS336 slide 43
\begin{frame}{How many heads, whats the model dim?}
% hand-typeset (_generator/handwritten/)
\begin{itemize}
  \item Some examples of this hyperparameter
\end{itemize}
\medskip
\begin{center}
\begin{tabular}{@{}lrrrr@{}}
\toprule
 & \textbf{Num heads} & \textbf{Head dim} & \textbf{Model dim} & \textbf{Ratio}\\
\midrule
GPT3 & 96 & 128 & 12288 & 1\\
T5 & 128 & 128 & 1024 & 16\\
T5 v1.1 & 64 & 64 & 4096 & 1\\
LaMDA & 128 & 128 & 8192 & 2\\
PaLM & 48 & 258 & 18432 & 1.48\\
LLaMA2 & 64 & 128 & 8192 & 1\\
\bottomrule
\end{tabular}
\end{center}
\medskip
Most models have ratios around 1 -- notable exceptions by some google models.
\end{frame}

% CS336 slide 44
\begin{frame}{Evidence for 1-1 ratio?}
\begin{itemize}
  \item There have been papers written against the 1-1 ratio [Bhojanapalli et al 2020]
  \item But we don't seem to be seeing significant `low rank bottlenecks' in practice..
\end{itemize}
\begin{center}
  \includegraphics[page=44,trim=210 95 210 115,clip,width=180pt]{figure/cs336-l3.pdf}\par
\end{center}
\end{frame}

% CS336 slide 45
\begin{frame}{Aspect ratios}
% hand-typeset (_generator/handwritten/)
\begin{itemize}
  \item Should my model be deep or wide? \emph{How} deep and how wide?
  \item Most models are surprisingly consistent on this one too!
\end{itemize}
\begin{center}
\small
\renewcommand{\arraystretch}{0.95}
\begin{tabular}{@{}llr@{}}
\toprule
 & \textbf{Model} & $\boldsymbol{d_{model}/n_{layer}}$\\
\midrule
 & BLOOM & 205\\
 & T5 v1.1 & 171\\
\rowcolor{blue!8} & PaLM (540B) & 156\\
\rowcolor{blue!8} \emph{Sweet spot?} & GPT3/OPT/Mistral/Qwen & 128\\
\rowcolor{blue!8} & LLaMA / LLaMA2 / Chinchila & 102\\
 & T5 (11B) & 43\\
 & GPT2 & 33\\
\bottomrule
\end{tabular}
\end{center}
\end{frame}

% CS336 slide 46
\begin{frame}{Considerations about aspect ratio}
% hand-typeset (_generator/handwritten/)
\begin{itemize}
  \item \textbf{Extremely deep models are harder to parallelize and have higher latency}
\end{itemize}
\begin{center}
\includegraphics[page=46,trim=76 208 39 117,clip,width=\linewidth]{figure/cs336-l3.pdf}\\[-1pt]
{\raggedleft\scriptsize [Tay et al 2021]\par}
\smallskip
\includegraphics[page=46,trim=189 10 257 209,clip,width=\linewidth,height=80pt,keepaspectratio]{figure/cs336-l3.pdf}
\end{center}
\end{frame}

% CS336 slide 47
\begin{frame}{Evidence on aspect ratio scaling}
% hand-typeset (_generator/handwritten/)
\begin{columns}[T]
\begin{column}{0.5\textwidth}
\centering
\includegraphics[page=47,trim=69 104 377 93,clip,width=\linewidth]{figure/cs336-l3.pdf}\\
{\scriptsize [Kaplan et al 2020]\par}
\end{column}
\begin{column}{0.48\textwidth}
\centering
\includegraphics[page=47,trim=392 67 71 77,clip,width=\linewidth,height=150pt,keepaspectratio]{figure/cs336-l3.pdf}\\
{\scriptsize [Tay et al 2021]\par}
\end{column}
\end{columns}
\end{frame}

% CS336 slide 48
\begin{frame}{What are typical vocabulary sizes?}
% hand-typeset (_generator/handwritten/)
\begin{columns}[T]
\begin{column}{0.48\textwidth}
\centering
Monolingual models -- 30-50k vocab

\medskip
\small
\begin{tabular}{@{}lr@{}}
\toprule
\textbf{Model} & \textbf{Token count}\\
\midrule
Original transformer & 37000\\
GPT & 40257\\
GPT2/3 & 50257\\
T5/T5v1.1 & 32128\\
LLaMA & 32000\\
\bottomrule
\end{tabular}
\end{column}
\begin{column}{0.48\textwidth}
\centering
Multilingual / production systems 100-250k

\medskip
\small
\begin{tabular}{@{}lr@{}}
\toprule
\textbf{Model} & \textbf{Token count}\\
\midrule
mT5 & 250000\\
PaLM & 256000\\
GPT4 & 100276\\
Command A & 255000\\
DeepSeek & 100000\\
Qwen 15B & 152064\\
Yi & 64000\\
\bottomrule
\end{tabular}
\end{column}
\end{columns}
\medskip
{\centering Monolingual vocabs don't need to be huge, but multilingual ones do\par}
\end{frame}

% CS336 slide 49
\begin{frame}{Dropout and other regularization}
\begin{itemize}
  \item Do we need regularization during pretraining?
  \item Arguments against:
  \begin{itemize}
    \item There is a lot of data (trillions of tokens), more than parameters.
    \item SGD only does a single pass on a corpus (hard to memorize)
  \end{itemize}
  \item This is all quite reasonable.. but what do people do in practice?
\end{itemize}
\end{frame}

% CS336 slide 50
\begin{frame}{Dropout and weight decay in practice}
% hand-typeset (_generator/handwritten/)
\begin{columns}[T]
\begin{column}{0.58\textwidth}
\footnotesize
\renewcommand{\arraystretch}{0.95}
\begin{tabular}{@{}lll@{}}
\toprule
\textbf{Model} & \textbf{Dropout*} & \textbf{Weight decay}\\
\midrule
Original transformer & 0.1 & 0\\
GPT2 & 0.1 & 0.1\\
T5 & 0.1 & 0\\
GPT3 & 0.1 & 0.1\\
T5 v1.1 & 0 & 0\\
PaLM & 0 & (variable)\\
OPT & 0.1 & 0.1\\
LLaMA & 0 & 0.1\\
Qwen 14B & 0.1 & 0.1\\
\bottomrule
\end{tabular}
\end{column}
\begin{column}{0.4\textwidth}
\vspace{2.5em}
\begin{itemize}
  \item Many older models used dropout during pretraining
  \item Newer models (except Qwen) rely only on weight decay
\end{itemize}
\end{column}
\end{columns}
\vfill
{\scriptsize * Most of the times papers just don't discuss dropout. On open models, this closely matches not doing dropout. This may not be true of closed models.\par}
\end{frame}

% CS336 slide 51
\begin{frame}{Why weight decay LLMs?}
% hand-typeset (_generator/handwritten/)
[Andriushchenko et al 2023] has interesting observations about LLM weight decay

\medskip
\begin{columns}[T]
\begin{column}{0.33\textwidth}
\centering
\includegraphics[page=51,trim=61 94 463 127,clip,width=\linewidth]{figure/cs336-l3.pdf}\\[2pt]
{\small It's not to control overfitting\par}
\end{column}
\begin{column}{0.65\textwidth}
\centering
\includegraphics[page=51,trim=275 106 19 127,clip,width=\linewidth]{figure/cs336-l3.pdf}\\[2pt]
{\small Weight decay interacts with learning rates (cosine schedule)\par}
\end{column}
\end{columns}
\end{frame}

% CS336 slide 52
\begin{frame}{Summary: hyperparameters}
% hand-typeset (_generator/handwritten/)
\makebox[\textwidth][c]{%
\begin{minipage}{400pt}
\begin{columns}[T,onlytextwidth]
\begin{column}{0.58\textwidth}
\footnotesize
\setlength{\itemsep}{0pt}
\textbf{Feedforward}
\begin{itemize}
  \item Factor-of-4 rule of thumb (8/3 for GLUs) is standard (with some evidence)
\end{itemize}
\textbf{Head dim}
\begin{itemize}
  \item Head dim*Num head = D model is standard -- but low to no validation
\end{itemize}
\textbf{Aspect ratio}
\begin{itemize}
  \item Wide range of `good' values (100-200). Systems concerns dictate the value
\end{itemize}
\textbf{Regularization}
\begin{itemize}
  \item You still `regularize' LMs but its effects are primarily on optimization dynamics
\end{itemize}
\end{column}
\begin{column}{0.4\textwidth}
\centering
\includegraphics[page=52,trim=395 14 35 44,clip,width=\linewidth,height=170pt,keepaspectratio]{figure/cs336-l3.pdf}
\end{column}
\end{columns}
\end{minipage}}
\end{frame}

% FRAMES-END

%%%%%%%%%%%%%%%%%%%%%%%%%%%%%%%%%%%%%%%%%%%%%%%%%%%%%%%%%%%%%%%%%%
\section{Stability Tricks}
% Source: CS336 Spring 2025, lecture 3 (architecture), slides 53-57
%%%%%%%%%%%%%%%%%%%%%%%%%%%%%%%%%%%%%%%%%%%%%%%%%%%%%%%%%%%%%%%%%%

% FRAMES-BEGIN l3 53 57
% CS336 slide 53
\begin{frame}{Stability tricks}
\begin{itemize}
  \item Recently, lots of attention on stable training
  \begin{itemize}
    \item Don't train models that look like the blue curve!
  \end{itemize}
\end{itemize}
\begin{center}
  \includegraphics[page=53,trim=163 70 163 102,clip,width=205pt]{figure/cs336-l3.pdf}\par
\end{center}
\end{frame}

% CS336 slide 54
\begin{frame}{Where do the issues arise? Beware of softmaxes!}
% hand-typeset (_generator/handwritten/)
\textbf{Softmaxes} -- can be ill-behaved due to exponentials / division by zero

\medskip
\begin{center}
\includegraphics[page=54,trim=167 34 220 124,clip,width=\linewidth,height=140pt,keepaspectratio]{figure/cs336-l3.pdf}
\end{center}
\end{frame}

% CS336 slide 55
\begin{frame}{Output softmax stability -- the `z-loss'}
% hand-typeset (_generator/handwritten/)
Recall the softmax calculation
\begin{center}
\includegraphics[page=55,trim=97 193 379 100,clip,width=\linewidth,height=62pt,keepaspectratio]{figure/cs336-l3.pdf}\hspace{10pt}\includegraphics[page=55,trim=366 220 66 100,clip,width=\linewidth,height=46pt,keepaspectratio]{figure/cs336-l3.pdf}\\[-1pt]
{\raggedleft\scriptsize [From Devlin 2014]\par}
\end{center}
This is useful for stability! PaLM pioneered this `z loss' trick.
\begin{center}
\includegraphics[page=55,trim=169 90 133 282,clip,width=0.85\linewidth]{figure/cs336-l3.pdf}
\end{center}
\textbf{Other examples:} Baichuan 2 (2023), DCLM (2024), OLMo 2 (2025)
\end{frame}

% CS336 slide 56
\begin{frame}{Attention softmax stability -- the `QK norm'}
% hand-typeset (_generator/handwritten/)
\begin{center}
\begin{tikzpicture}
\node[anchor=south west, inner sep=0] (img) {\includegraphics[page=56,trim=126 86 126 81,clip,width=\linewidth,height=118pt,keepaspectratio]{figure/cs336-l3.pdf}};
\begin{scope}[x={(img.south east)}, y={(img.north west)}]
\fill[white] (0.000,0.080) rectangle (0.212,0.214);
\fill[white] (0.842,0.080) rectangle (1.013,0.214);
\end{scope}
\end{tikzpicture}
\end{center}
The query and keys are Layer (RMS) normed before going into the softmax operation.

\medskip
\textbf{Other examples: DCLM, OLMo2, Gemma 2}\\
{\footnotesize Originally from vision and multimodal models [Dehgani 2023, Idefcs, Chameleon]\par}
\end{frame}

% CS336 slide 57
\begin{frame}{Logit soft-capping.}
\begin{itemize}
  \item Soft-capping the logits to some maximum value via Tanh
  \item Prevents logits from blowing up, but also might have perf issues?
\end{itemize}
\begin{center}
  \includegraphics[page=57,trim=249 141 212 106,clip,width=151pt]{figure/cs336-l3.pdf}\par
  \includegraphics[page=57,trim=202 38 166 319,clip,width=206pt]{figure/cs336-l3.pdf}\par
\end{center}
\end{frame}

% FRAMES-END

%%%%%%%%%%%%%%%%%%%%%%%%%%%%%%%%%%%%%%%%%%%%%%%%%%%%%%%%%%%%%%%%%%
\section{Attention Heads: MQA, GQA, Sliding Windows}
% Source: CS336 Spring 2025, lecture 3 (architecture), slides 58-58
%%%%%%%%%%%%%%%%%%%%%%%%%%%%%%%%%%%%%%%%%%%%%%%%%%%%%%%%%%%%%%%%%%

% FRAMES-BEGIN l3 58 58
% CS336 slide 58
\begin{frame}{Attention heads}
\begin{itemize}
  \item Most models don't touch the attention heads much at all with a few minor exceptions..
  \item GQA / MQA : Saving inference costs by reducing the number of heads
  \item Sparse or sliding window attention (GPT4/Mistral): restricting the attention pattern
  \begin{itemize}
    \item to reduce compute cost
  \end{itemize}
  \item Exotic SSM stuff (Jamba, Falcon 3, etc): not covered (sorry!)
\end{itemize}
\end{frame}

% FRAMES-END

% [O] original bridge frame (draft for the user to rewrite); listed in order.txt
\begin{frame}{The KV Cache}
\begin{itemize}
  \item Autoregressive generation (session 8, decoding) produces one token at a time; each new token attends to all previous tokens.
  \item The keys and values of the previous tokens do not change, so they are stored instead of recomputed: the \textbf{KV cache}.
  \item Its size grows with sequence length $\times$ number of layers $\times$ number of heads $\times$ head dimension (per sequence in the batch).
  \item At inference time, reading this cache from memory --- not the arithmetic --- is often the bottleneck. The following attention variants reduce it.
\end{itemize}
\end{frame}

% FRAMES-BEGIN l3 59 68
% CS336 slide 59
\begin{frame}{GQA/MQA -- Reducing attention head cost}
% hand-typeset (_generator/handwritten/)
\small
\textbf{Let's think about the compute involved for attention}
\begin{center}
\includegraphics[page=59,trim=153 127 201 109,clip,width=\linewidth,height=80pt,keepaspectratio]{figure/cs336-l3.pdf}
\end{center}
\textbf{Total arithmetic operations} $(bnd^2)$, \textbf{total memory accesses}
\[ (\overset{\text{\scriptsize X}}{bnd} + \overset{\text{\scriptsize softmax}}{bhn^2} + \overset{\text{\scriptsize projection}}{d^2}) \]

Arithmetic intensity is high $O\left(\left(\frac{1}{k} + \frac{1}{bn}\right)^{-1}\right)$ -- we can keep our GPUs running
\end{frame}

% CS336 slide 60
\begin{frame}[shrink=8]{GQA/MQA -- Reducing attention head cost}
\begin{itemize}
  \item What about the incremental case when we generate text?
  \item Key difference: can't parallelize the generation process -- needs to be step by step
  \item In this case -- we need to incrementaly re-compute/update attention via the `KV cache'
\end{itemize}
\begin{center}
  \includegraphics[page=60,trim=191 22 195 194,clip,width=185pt]{figure/cs336-l3.pdf}\par
\end{center}
{\scriptsize [Animation from https://medium.com/@joaolages/kv-caching-explained-276520203249]\par}
\end{frame}

% CS336 slide 61
\begin{frame}{GQA/MQA -- Reducing attention head cost}
% hand-typeset (_generator/handwritten/)
What's the incremental arithmetic intensity?

\bigskip
\textbf{Total arithmetic operations} $(bnd^2)$, \textbf{total memory accesses}
\[ (\overset{\text{\scriptsize X}}{bn^2 d} + \overset{\text{\scriptsize projection}}{nd^2}) \]

\bigskip
Arithmetic intensity is not good $O\left(\left(\frac{n}{d} + \frac{1}{b}\right)^{-1}\right)$ -- need large batches + short seq length (n) or big model dimensions (d)

\bigskip
Is there some way around this? The n/d term is difficult to reduce.
\end{frame}

% CS336 slide 62
\begin{frame}{MQA -- just have fewer key dimensions.}
% hand-typeset (_generator/handwritten/)
\small
\textbf{Key idea} -- have multiple queries, but just one dimension for keys and values
\begin{center}
\includegraphics[page=62,trim=181 91 181 112,clip,width=\linewidth,height=82pt,keepaspectratio]{figure/cs336-l3.pdf}
\end{center}
We have much fewer items to move in and out of memory (KV Cache)

\smallskip
\textbf{Total memory access} $(bnd + bn^2 k + nd^2)$, \textbf{Arithmetic intensity} $O\big(\big(\textcolor{red}{\tfrac{1}{d}} + \tfrac{n}{\textcolor{red}{d}h} + \tfrac{1}{b}\big)^{-1}\big)$

{\raggedleft\tiny [figure from \url{https://blog.fireworks.ai/multi-query-attention-is-all-you-need-db072e758055}]\par}
\end{frame}

% CS336 slide 63
\begin{frame}{Recent extension -- GQA}
\begin{itemize}
  \item Don't go all the way to one dimension of KV -- have fewer dims
  \item Simple knob to control expressiveness (key-query ratio) and inference efficiency
\end{itemize}
\begin{center}
  \includegraphics[page=63,trim=117 119 130 125,clip,width=284pt]{figure/cs336-l3.pdf}\par
\end{center}
\end{frame}

% CS336 slide 64
\begin{frame}{Does MQA hurt? Sometimes..}
% hand-typeset (_generator/handwritten/)
\begin{columns}[T]
\begin{column}{0.49\textwidth}
\centering
Small PPL hit w/ MQA [Shazeer 2019]\\[2pt]
\includegraphics[page=64,trim=98 116 389 156,clip,width=\linewidth]{figure/cs336-l3.pdf}
\end{column}
\begin{column}{0.49\textwidth}
\centering
Low/no hit w/ GQA [Ainslie 2023]\\[2pt]
\includegraphics[page=64,trim=437 154 117 115,clip,width=\linewidth,height=70pt,keepaspectratio]{figure/cs336-l3.pdf}\\
\includegraphics[page=64,trim=405 6 77 251,clip,width=\linewidth,height=75pt,keepaspectratio]{figure/cs336-l3.pdf}
\end{column}
\end{columns}
\end{frame}

% CS336 slide 65
\begin{frame}{Sparse / sliding window attention}
\begin{itemize}
  \item Attending to the entire context can be expensive (quadratic).
  \item Build sparse / structured attention that trades off expressiveness vs runtime (GPT3)
\end{itemize}
\begin{center}
  \includegraphics[page=65,trim=119 39 162 146,clip,width=241pt]{figure/cs336-l3.pdf}\par
\end{center}
{\scriptsize [Child et al 2019]\par}
\end{frame}

% CS336 slide 66
\begin{frame}{Sliding window attention}
\begin{itemize}
  \item Another variation on this idea -- sliding window attention
  \item Just use the main part of the strided pattern -- let depth extend effective context (Mistral)
\end{itemize}
\begin{center}
  \includegraphics[page=66,trim=90 93 118 106,clip,width=301pt]{figure/cs336-l3.pdf}\par
\end{center}
\end{frame}

% CS336 slide 67
\begin{frame}{Current standard trick -- interleave `full' and `LR' attention}
\begin{itemize}
  \item From Cohere Command A -- Every 4th layer is a full attention
  \item Long-range info via NoPE, short-range info via RoPE + SWA.
  \item Other models -- LLaMA 4, Gemma does SWA+Full RoPE.
\end{itemize}
\begin{center}
  \includegraphics[page=67,trim=178 104 178 104,clip,width=193pt]{figure/cs336-l3.pdf}\par
\end{center}
\end{frame}

% CS336 slide 68
\begin{frame}{Recap, conclusion, etc.}
\begin{itemize}
  \item Many aspects (arch, hparams) of transformers are in common across the big LMs
  \begin{itemize}
    \item Major differences? Position embeddings, activations, tokenization
  \end{itemize}
\end{itemize}
\begin{center}
  \includegraphics[page=68,trim=35 92 21 106,clip,width=330pt]{figure/cs336-l3.pdf}\par
\end{center}
\end{frame}

% FRAMES-END

%%%%%%%%%%%%%%%%%%%%%%%%%%%%%%%%%%%%%%%%%%%%%%%%%%%%%%%%%%%%%%%%%%
\section{Mixture of Experts}
% Source: CS336 Spring 2025, lecture 4 (MoEs), slides 2-12
%%%%%%%%%%%%%%%%%%%%%%%%%%%%%%%%%%%%%%%%%%%%%%%%%%%%%%%%%%%%%%%%%%

% FRAMES-BEGIN l4 2 12
% CS336 slide 2
\begin{frame}{Mixture of experts}
\begin{center}
  \includegraphics[page=2,trim=30 24 12 68,clip,width=330pt]{figure/cs336-l4.pdf}\par
\end{center}
\end{frame}

% CS336 slide 3
\begin{frame}{What's a MoE?}
% hand-typeset (_generator/handwritten/)
\begin{center}
\includegraphics[page=3,trim=51 128 51 101,clip,width=\linewidth]{figure/cs336-l4.pdf}
\end{center}
{\raggedleft\scriptsize [Fedus et al 2022]\par}
\begin{itemize}
  \item Replace big feedforward with (many) big feedforward networks and a selector layer
  \begin{itemize}
    \item You can increase the \# experts without affecting FLOPs
  \end{itemize}
\end{itemize}
\end{frame}

% CS336 slide 4
\begin{frame}{Why are MoEs getting popular?}
% hand-typeset (_generator/handwritten/)
\begin{itemize}
  \item Same FLOP, more param does better
\end{itemize}
\begin{center}
\includegraphics[page=4,trim=79 59 66 112,clip,width=\linewidth,height=120pt,keepaspectratio]{figure/cs336-l4.pdf}
\end{center}
{\raggedleft\scriptsize [Fedus et al 2022]\par}
\end{frame}

% CS336 slide 5
\begin{frame}{Why are MoEs getting popular?}
% hand-typeset (_generator/handwritten/)
\begin{itemize}
  \item Faster to train MoEs
\end{itemize}
\begin{center}
\includegraphics[page=5,trim=61 36 349 121,clip,width=0.47\linewidth]{figure/cs336-l4.pdf}\hfill\includegraphics[page=5,trim=368 52 6 138,clip,width=0.51\linewidth]{figure/cs336-l4.pdf}
\end{center}
{\raggedleft\scriptsize [OlMoE]\par}
\end{frame}

% CS336 slide 6
\begin{frame}{Why are MoEs getting popular?}
\begin{itemize}
  \item Highly competitive vs dense equivalents
\end{itemize}
\begin{center}
  \includegraphics[page=6,trim=183 55 183 88,clip,width=185pt]{figure/cs336-l4.pdf}\par
\end{center}
\end{frame}

% CS336 slide 7
\begin{frame}{Why are MoEs getting popular?}
% hand-typeset (_generator/handwritten/)
\begin{itemize}
  \item Parallelizable to many devices
\end{itemize}
\begin{center}
\includegraphics[page=7,trim=181 14 200 101,clip,width=\linewidth,height=140pt,keepaspectratio]{figure/cs336-l4.pdf}
\end{center}
\end{frame}

% CS336 slide 8
\begin{frame}{Some MoE results -- from the west}
\begin{itemize}
  \item MoEs are most of the highest-performance open models, and are quite quick
\end{itemize}
\begin{center}
  \includegraphics[page=8,trim=116 82 116 82,clip,width=277pt]{figure/cs336-l4.pdf}\par
\end{center}
\end{frame}

% CS336 slide 9
\begin{frame}{Earlier MoE results from Chinese groups -- Qwen}
\begin{itemize}
  \item Chinese LLM companies are also doing quite a bit of MoE work on the smaller end
\end{itemize}
\begin{center}
  \includegraphics[page=9,trim=71 106 315 136,clip,width=140pt]{figure/cs336-l4.pdf}\par
  \includegraphics[page=9,trim=408 106 54 136,clip,width=109pt]{figure/cs336-l4.pdf}\par
\end{center}
\end{frame}

% CS336 slide 10
\begin{frame}{Earlier MoE results from Chinese groups - DeepSeek}
\begin{itemize}
  \item There's also some good recent ablation work on MoEs showing they're generally good
\end{itemize}
\begin{center}
  \includegraphics[page=10,trim=193 20 193 107,clip,width=165pt]{figure/cs336-l4.pdf}\par
\end{center}
\end{frame}

% CS336 slide 11
\begin{frame}{Recent MoE results -- DeepSeek v3}
\begin{center}
  \includegraphics[page=11,trim=119 30 82 67,clip,width=277pt]{figure/cs336-l4.pdf}\par
\end{center}
\end{frame}

% CS336 slide 12
\begin{frame}{Why haven't MoEs been more popular?}
% hand-typeset (_generator/handwritten/)
\begin{itemize}
  \item Infrastructure is complex / advantages on multi node
\end{itemize}
\begin{center}
\includegraphics[page=12,trim=113 193 113 104,clip,width=\linewidth,height=46pt,keepaspectratio]{figure/cs336-l4.pdf}\\[-1pt]
\end{center}
\vspace{-8pt}
{\raggedleft\scriptsize [Fedus et al 2022]\par}
\begin{itemize}
  \item Training objectives are somewhat heuristic (and sometimes unstable)
\end{itemize}
\begin{center}
\includegraphics[page=12,trim=194 1 221 263,clip,width=\linewidth,height=60pt,keepaspectratio]{figure/cs336-l4.pdf}
\end{center}
\vspace{-8pt}
{\raggedleft\scriptsize [Zoph et al 2022]\par}
\end{frame}

% FRAMES-END

%%%%%%%%%%%%%%%%%%%%%%%%%%%%%%%%%%%%%%%%%%%%%%%%%%%%%%%%%%%%%%%%%%
\section{MoE Routing}
% Source: CS336 Spring 2025, lecture 4 (MoEs), slides 13-23
%%%%%%%%%%%%%%%%%%%%%%%%%%%%%%%%%%%%%%%%%%%%%%%%%%%%%%%%%%%%%%%%%%

% FRAMES-BEGIN l4 13 23
% CS336 slide 13
\begin{frame}{What MoEs generally look like}
% hand-typeset (_generator/handwritten/)
\begin{columns}[T]
\begin{column}{0.58\textwidth}
\centering
Typical: replace MLP with MoE layer\\[6pt]
\includegraphics[page=13,trim=46 68 357 161,clip,width=\linewidth]{figure/cs336-l4.pdf}
\end{column}
\begin{column}{0.4\textwidth}
\centering
Less common: MoE for attention heads\\[6pt]
\includegraphics[page=13,trim=476 44 95 149,clip,width=\linewidth,height=125pt,keepaspectratio]{figure/cs336-l4.pdf}\\
{\scriptsize [ModuleFormer, JetMoE]\par}
\end{column}
\end{columns}
\end{frame}

% CS336 slide 14
\begin{frame}{MoE -- what varies?}
\begin{itemize}
  \item Routing function
  \item Expert sizes
  \item Training objectives
\end{itemize}
\end{frame}

% CS336 slide 15
\begin{frame}{Routing function - overview}
% hand-typeset (_generator/handwritten/)
\begin{itemize}
  \item Many of the routing algorithms boil down to `choose top k'
\end{itemize}
\begin{center}
\includegraphics[page=15,trim=121 76 121 118,clip,width=0.95\linewidth]{figure/cs336-l4.pdf}\\[2pt]
\makebox[0.95\linewidth]{\small
\makebox[0.33\linewidth]{\parbox{0.3\linewidth}{\centering Token chooses expert}}%
\makebox[0.33\linewidth]{\parbox{0.3\linewidth}{\centering Expert chooses token}}%
\makebox[0.33\linewidth]{\parbox{0.3\linewidth}{\centering Global routing via optimization}}}
\end{center}
{\raggedleft\scriptsize [Fedus et al 2022]\par}
\end{frame}

% CS336 slide 16
\begin{frame}{Routing type}
\begin{itemize}
  \item Almost all the MoEs do a standard `token choice topk' routing. Some recent ablations
\end{itemize}
\begin{center}
  \includegraphics[page=16,trim=51 119 51 119,clip,width=330pt]{figure/cs336-l4.pdf}\par
\end{center}
\end{frame}

% CS336 slide 17
\begin{frame}{Common routing variants in detail}
% hand-typeset (_generator/handwritten/)
\small
\makebox[\textwidth][c]{%
\begin{tabular}{@{}>{\raggedright\arraybackslash}m{50pt}@{\hspace{8pt}}c@{\hspace{14pt}}>{\raggedright\arraybackslash}m{160pt}@{}}
Top-k & \includegraphics[page=17,trim=125 173 345 70,clip,width=170pt,height=64pt,keepaspectratio]{figure/cs336-l4.pdf}
 & Used in \emph{most} MoEs\par\smallskip
   Switch Transformer (k=1)\newline
   Gshard (k=2), Grok (2), Mixtral (2), Qwen (4), DBRX (4), DeepSeek (7)\\[2pt]
Hashing & \includegraphics[page=17,trim=125 -2 345 244,clip,width=170pt,height=64pt,keepaspectratio]{figure/cs336-l4.pdf}
 & Common baseline
\end{tabular}}

{\raggedleft\scriptsize [Fedus et al 2022]\par}
\end{frame}

% CS336 slide 18
\begin{frame}{Other routing methods}
% hand-typeset (_generator/handwritten/)
\renewcommand{\arraystretch}{1.2}
\makebox[\textwidth][c]{%
\begin{tabular}{@{}>{\raggedright\arraybackslash}m{70pt}@{\hspace{8pt}}c@{\hspace{10pt}}>{\raggedright\arraybackslash}m{170pt}@{}}
RL to learn routes & \includegraphics[page=18,trim=171 176 307 71,clip,width=\linewidth,height=75pt,keepaspectratio]{figure/cs336-l4.pdf} & Used in some of the earliest work\newline Bengio 2013, not common now\\
Solve a matching problem & \includegraphics[page=18,trim=171 8 307 240,clip,width=\linewidth,height=75pt,keepaspectratio]{figure/cs336-l4.pdf} & Linear assignment for routing\newline Used in various papers like Clark `22
\end{tabular}}

{\raggedleft\scriptsize [Fedus et al 2022]\par}
\end{frame}

% CS336 slide 19
\begin{frame}{Top-K routing in detail.}
\begin{itemize}
  \item Most papers do the old and classic top-k routing. How does this work?
  \begin{itemize}
    \item This is the DeepSeek (V1-2) router (Grok, Qwen do this too)
    \item Mixtral, DBRX, DeepSeek v3 softmaxes after the TopK
  \end{itemize}
  \item Gates selected by a logistic regressor
\end{itemize}
\begin{center}
  \includegraphics[page=19,trim=21 46 275 112,clip,width=153pt]{figure/cs336-l4.pdf}\par
\end{center}
{\scriptsize [Dai et al 2024]\par}
\end{frame}

% CS336 slide 20
\begin{frame}{Recent variations from DeepSeek and other Chinese LMs}
\begin{itemize}
  \item Smaller, larger number of experts + a few shared experts that are always on.
  \begin{itemize}
    \item (Used in DeepSeek / Qwen, originally from DeepSpeed MoE)
  \end{itemize}
\end{itemize}
\begin{center}
  \includegraphics[page=20,trim=117 86 81 64,clip,width=247pt]{figure/cs336-l4.pdf}\par
\end{center}
\end{frame}

% CS336 slide 21
\begin{frame}{Various ablations from the DeepSeek paper}
\begin{itemize}
  \item More experts, shared experts all seem to generally help
\end{itemize}
\begin{center}
  \includegraphics[page=21,trim=128 66 128 69,clip,width=236pt]{figure/cs336-l4.pdf}\par
\end{center}
\end{frame}

% CS336 slide 22
\begin{frame}{Ablations from OlMoE}
\begin{itemize}
  \item Gains from fine-grained experts, none from shared experts.
\end{itemize}
\begin{center}
  \includegraphics[page=22,trim=70 4 80 100,clip,width=260pt]{figure/cs336-l4.pdf}\par
\end{center}
\end{frame}

% CS336 slide 23
\begin{frame}{Expert routing setups for recent MoEs}
% hand-typeset (_generator/handwritten/)
\begin{center}
\footnotesize
\renewcommand{\arraystretch}{0.92}
\begin{tabular}{@{}lrrrc@{}}
\toprule
\textbf{Model} & \textbf{Routed} & \textbf{Active} & \textbf{Shared} & \textbf{Fine-grained ratio}\\
\midrule
GShard & 2048 & 2 & 0 & \\
Switch Transformer & 64 & 1 & 0 & \\
ST-MOE & 64 & 2 & 0 & \\
Mixtral & 8 & 2 & 0 & \\
DBRX & 16 & 4 & 0 & \\
Grok & 8 & 2 & 0 & \\
DeepSeek v1 & 64 & 6 & 2 & 1/4\\
Qwen 1.5 & 60 & 4 & 4 & 1/8\\\midrule
DeepSeek v3 & 256 & 8 & 1 & 1/14\\
OlMoE & 64 & 8 & 0 & 1/8\\
MiniMax & 32 & 2 & 0 & $\sim$1/4\\
Llama 4 (maverick) & 128 & 1 & 1 & 1/2\\
\bottomrule
\end{tabular}
\end{center}
\end{frame}

% FRAMES-END

%%%%%%%%%%%%%%%%%%%%%%%%%%%%%%%%%%%%%%%%%%%%%%%%%%%%%%%%%%%%%%%%%%
\section{Training MoEs}
% Source: CS336 Spring 2025, lecture 4 (MoEs), slides 24-37
%%%%%%%%%%%%%%%%%%%%%%%%%%%%%%%%%%%%%%%%%%%%%%%%%%%%%%%%%%%%%%%%%%

% FRAMES-BEGIN l4 24 37
% CS336 slide 24
\begin{frame}{How do we train MoEs?}
\begin{itemize}
  \item Major challenge: we need sparsity for training-time efficiency\ldots{}
  \begin{itemize}
    \item But sparse gating decisions are not differentiable!
  \end{itemize}
  \item Solutions?
  \item 1. Reinforcment learning to optimize gating policies 2. Stochastic perturbations 3. Heuristic `balancing' losses.
  \begin{itemize}
    \item Guess which one people use in practice?
  \end{itemize}
\end{itemize}
\end{frame}

% CS336 slide 25
\begin{frame}{RL for MoEs}
\begin{itemize}
  \item RL via REINFORCE does work, but not so much better that it's a clear win
  \begin{itemize}
    \item RL is the `right solution' but gradient variances and complexity
    \item means it's not widely used
  \end{itemize}
\end{itemize}
\begin{center}
  \includegraphics[page=25,trim=55 122 35 137,clip,width=330pt]{figure/cs336-l4.pdf}\par
\end{center}
{\scriptsize (REINFORCE baseline approach, Clark et al 2020)\par}
\end{frame}

% CS336 slide 26
\begin{frame}{Stochastic approximations}
\begin{itemize}
  \item From Shazeer et al 2017 -- routing decisions are stochastic with gaussian perturbations.
  \item 1. This naturally leads to experts that are a bit more robust. 2. The softmax means that the model learns how to rank K experts
\end{itemize}
\begin{center}
  \includegraphics[page=26,trim=106 137 106 90,clip,width=305pt]{figure/cs336-l4.pdf}\par
\end{center}
\end{frame}

% CS336 slide 27
\begin{frame}{Stochastic approximations}
\begin{itemize}
  \item Stochastic jitter in Fedus et al 2022. This does a uniform multiplicative perturbation for the
  \item same goal of getting less brittle experts. This was later removed in Zoph et al 2022
\end{itemize}
\begin{center}
  \includegraphics[page=27,trim=98 187 61 86,clip,width=315pt]{figure/cs336-l4.pdf}\par
  \includegraphics[page=27,trim=212 21 212 301,clip,width=167pt]{figure/cs336-l4.pdf}\par
\end{center}
\end{frame}

% CS336 slide 28
\begin{frame}{Heuristic balancing losses}
% hand-typeset (_generator/handwritten/)
\begin{itemize}
  \item Another key issue -- systems efficiency requires that we use experts evenly..
\end{itemize}
\begin{center}
\includegraphics[page=28,trim=121 87 159 103,clip,width=\linewidth,height=108pt,keepaspectratio]{figure/cs336-l4.pdf}
\end{center}
\vspace{-8pt}
{\raggedleft\tiny From the Switch Transformer [Fedus et al 2022]\par}
{\centering\small The derivative with respect to $p_i(x)$ is $\frac{\alpha N}{T^2} \sum \mathbb{1}_{\mathit{argmax}\, p(x) = i}$,\\ so more frequent use = stronger downweighting\par}
\end{frame}

% CS336 slide 29
\begin{frame}{Example from deepseek (v1-2)}
\begin{itemize}
  \item Per-expert balancing -- same as the switch transformer
  \item Per-device balancing -- the objective above, but aggregated by device.
\end{itemize}
\begin{center}
  \includegraphics[page=29,trim=114 182 186 109,clip,width=230pt]{figure/cs336-l4.pdf}\par
  \includegraphics[page=29,trim=111 17 183 281,clip,width=233pt]{figure/cs336-l4.pdf}\par
\end{center}
\end{frame}

% CS336 slide 30
\begin{frame}{DeepSeek v3 variation -- per-expert biases}
% hand-typeset (_generator/handwritten/)
\begin{itemize}
  \item Set up a per-expert bias (making it more likely to get tokens) and use online learning
\end{itemize}
\begin{center}
\includegraphics[page=30,trim=173 237 173 104,clip,width=\linewidth,height=22pt,keepaspectratio]{figure/cs336-l4.pdf}
\end{center}
\begin{itemize}
  \item They call this `\textbf{auxiliary loss free balancing}'
\end{itemize}
\begin{center}
\includegraphics[page=30,trim=142 27 142 198,clip,width=\linewidth,height=68pt,keepaspectratio]{figure/cs336-l4.pdf}
\end{center}
{\small (but the approach is not fully aux loss free..)\par}
\end{frame}

% CS336 slide 31
\begin{frame}{What happens when removing load balancing losses?}
\begin{center}
  \includegraphics[page=31,trim=158 14 158 74,clip,width=210pt]{figure/cs336-l4.pdf}\par
\end{center}
\end{frame}

% CS336 slide 32
\begin{frame}{Training MoEs -- the systems side}
% hand-typeset (_generator/handwritten/)
\begin{columns}[T]
\begin{column}{0.36\textwidth}
\centering
MoEs parallelize nicely -- Each FFN can fit in a device\\[4pt]
\includegraphics[page=32,trim=63 40 475 142,clip,width=\linewidth,height=125pt,keepaspectratio]{figure/cs336-l4.pdf}
\end{column}
\begin{column}{0.62\textwidth}
\centering
Enables additional kinds of parallelism\\[4pt]
\includegraphics[page=32,trim=307 38 21 139,clip,width=\linewidth]{figure/cs336-l4.pdf}
\end{column}
\end{columns}
\end{frame}

% CS336 slide 33
\begin{frame}{Training MoEs -- the systems side}
\begin{itemize}
  \item MoE routing allows for parallelism, but also some complexities
  \item Modern libraries like MegaBlocks (used in many open MoEs) use smarter sparse MMs
\end{itemize}
\begin{center}
  \includegraphics[page=33,trim=56 112 51 125,clip,width=330pt]{figure/cs336-l4.pdf}\par
\end{center}
\end{frame}

% CS336 slide 34
\begin{frame}{Fun side issue -- stochasticity of MoE models}
\begin{itemize}
  \item There was speculation that GPT-4's stochasticity was due to MoE..
  \item Why would a MoE have additional randomness?
  \begin{itemize}
    \item Token dropping from routing happens at a batch level -- this means that
    \item other people's queries can drop your token!
  \end{itemize}
\end{itemize}
\begin{center}
  \includegraphics[page=34,trim=101 79 101 155,clip,width=270pt]{figure/cs336-l4.pdf}\par
\end{center}
\end{frame}

% CS336 slide 35
\begin{frame}{Issues with MoEs - stability}
% hand-typeset (_generator/handwritten/)
\begin{center}
\includegraphics[page=35,trim=55 214 509 77,clip,width=0.25\linewidth]{figure/cs336-l4.pdf}\hspace{6pt}\includegraphics[page=35,trim=221 221 19 84,clip,width=0.72\linewidth]{figure/cs336-l4.pdf}
\end{center}
\vspace{-8pt}
{\raggedleft\scriptsize [Zoph 2022]\par}
\begin{itemize}
  \item \textbf{Solution:} Use Float 32 just for the expert router (sometimes with an aux z-loss)
\end{itemize}
\begin{center}
\includegraphics[page=35,trim=72 64 143 271,clip,width=0.85\linewidth]{figure/cs336-l4.pdf}
\end{center}
\end{frame}

% CS336 slide 36
\begin{frame}{Z-loss stability for the router}
\begin{itemize}
  \item What happens when we remove the z-loss?
\end{itemize}
\begin{center}
  \includegraphics[page=36,trim=88 100 88 100,clip,width=327pt]{figure/cs336-l4.pdf}\par
\end{center}
\end{frame}

% CS336 slide 37
\begin{frame}{Issues with MoEs -- fine-tuning}
% hand-typeset (_generator/handwritten/)
\begin{center}
Sparse MoEs can overfit on smaller fine-tuning data\\[2pt]
\includegraphics[page=37,trim=266 176 282 97,clip,width=\linewidth,height=70pt,keepaspectratio]{figure/cs336-l4.pdf}
\end{center}
\begin{columns}[T]
\begin{column}{0.4\textwidth}
\centering
{\small Zoph et al solution -- finetune non-MoE MLPs\par}
\smallskip
\includegraphics[page=37,trim=77 11 493 267,clip,width=\linewidth,height=62pt,keepaspectratio]{figure/cs336-l4.pdf}
\end{column}
\begin{column}{0.58\textwidth}
\centering
{\small DeepSeek solution -- use lots of data 1.4M SFT\par}
\smallskip
\includegraphics[page=37,trim=321 48 25 294,clip,width=\linewidth]{figure/cs336-l4.pdf}
\end{column}
\end{columns}
\end{frame}

% FRAMES-END

%%%%%%%%%%%%%%%%%%%%%%%%%%%%%%%%%%%%%%%%%%%%%%%%%%%%%%%%%%%%%%%%%%
\section{Upcycling and DeepSeek MoE}
% Source: CS336 Spring 2025, lecture 4 (MoEs), slides 38-47
%%%%%%%%%%%%%%%%%%%%%%%%%%%%%%%%%%%%%%%%%%%%%%%%%%%%%%%%%%%%%%%%%%

% FRAMES-BEGIN l4 38 47
% CS336 slide 38
\begin{frame}{Other training methods - upcycling}
\begin{itemize}
  \item Can we use a pre-trained LM to initialize a MoE?
\end{itemize}
\begin{center}
  \includegraphics[page=38,trim=459 110 30 101,clip,width=84pt]{figure/cs336-l4.pdf}\par
  \includegraphics[page=38,trim=33 117 271 104,clip,width=150pt]{figure/cs336-l4.pdf}\par
\end{center}
\end{frame}

% CS336 slide 39
\begin{frame}{Upcycling example - MiniCPM}
\begin{itemize}
  \item Uses the MiniCPM model (topk=2, 8 experts, \textasciitilde{} 4B active params).
  \item Simple MoE, shows gains from the base model with \textasciitilde{} 520B tokens for training
\end{itemize}
\begin{center}
  \includegraphics[page=39,trim=148 132 115 126,clip,width=274pt]{figure/cs336-l4.pdf}\par
\end{center}
\end{frame}

% CS336 slide 40
\begin{frame}{Upcycling example -- Qwen MoE}
\begin{itemize}
  \item Qwen MoE -- Initialized from the Qwen 1.8B model top-k=4, 60 experts w/ 4 shared.
  \item Similar architecture / setup to DeepSeekMoE, but one of the first (confirmed) upcycling successes
\end{itemize}
\begin{center}
  \includegraphics[page=40,trim=117 114 109 135,clip,width=297pt]{figure/cs336-l4.pdf}\par
\end{center}
\end{frame}

% CS336 slide 41
\begin{frame}{DeepSeek MoE v1-v2-v3}
% hand-typeset (_generator/handwritten/)
To wrap up, we'll walk through the DeepSeek MoE architecture.

\smallskip
\textbf{V1 (16B -- 2.8 active):}
\begin{center}
{\footnotesize\textbf{Shared (2) + Fine-grained (64/4) experts}\par}
\includegraphics[page=41,trim=207 149 217 145,clip,width=\linewidth,height=56pt,keepaspectratio]{figure/cs336-l4.pdf}
\end{center}
\begin{columns}[T]
\begin{column}{0.48\textwidth}
\centering
{\footnotesize\textbf{Standard, top-k routing}\par}
\includegraphics[page=41,trim=90 40 467 291,clip,width=\linewidth,height=38pt,keepaspectratio]{figure/cs336-l4.pdf}
\end{column}
\begin{column}{0.48\textwidth}
\centering
{\footnotesize\textbf{Standard Aux-loss balancing} (Expert + Device)\par}
\includegraphics[page=41,trim=385 34 193 296,clip,width=\linewidth,height=38pt,keepaspectratio]{figure/cs336-l4.pdf}
\end{column}
\end{columns}
\end{frame}

% CS336 slide 42
\begin{frame}{DeepSeek MoE v2}
% hand-typeset (_generator/handwritten/)
\textbf{V2 (236B -- 21 active):}
\begin{center}
{\footnotesize\textbf{Shared (2) + Fine-grained (160/10) experts, 6 active}\par}
\includegraphics[page=42,trim=211 174 211 122,clip,width=\linewidth,height=55pt,keepaspectratio]{figure/cs336-l4.pdf}
\end{center}
\textbf{\footnotesize New things:}
\begin{columns}[T]
\begin{column}{0.44\textwidth}
\centering
{\footnotesize Top-M device routing\par}
\includegraphics[page=42,trim=59 56 415 298,clip,width=\linewidth]{figure/cs336-l4.pdf}
\end{column}
\begin{column}{0.54\textwidth}
\centering
{\footnotesize Communication balancing loss -- balancing both communication in \emph{and} out\par}
\includegraphics[page=42,trim=375 5 109 294,clip,width=\linewidth,height=52pt,keepaspectratio]{figure/cs336-l4.pdf}
\end{column}
\end{columns}
\end{frame}

% CS336 slide 43
\begin{frame}{DeepSeek MoE v3}
% hand-typeset (_generator/handwritten/)
\textbf{V3 (671B -- 37 active):}
\begin{center}
{\footnotesize\textbf{Shared (1) + Fine-grained (258) experts, 8 active}\par}
\includegraphics[page=43,trim=211 176 211 122,clip,width=\linewidth,height=55pt,keepaspectratio]{figure/cs336-l4.pdf}
\end{center}
\textbf{\small New things}
\begin{columns}[T]
\begin{column}{0.48\textwidth}
\centering
{\footnotesize Sigmoid+Softmax topK + topM\par}
\includegraphics[page=43,trim=110 14 401 268,clip,width=\linewidth,height=55pt,keepaspectratio]{figure/cs336-l4.pdf}
\end{column}
\begin{column}{0.48\textwidth}
\centering
{\footnotesize Aux-loss-free + seq-wise aux\par}
\vspace{10pt}
\includegraphics[page=43,trim=415 54 81 312,clip,width=\linewidth]{figure/cs336-l4.pdf}
\end{column}
\end{columns}
\end{frame}

% CS336 slide 44
\begin{frame}{Bonus: What else do you need to make DeepSeek MoE v3?}
\begin{itemize}
  \item MLA : Multihead, latent attention
  \item Basic idea: express the Q, K, V as functions of a lower-dim, `latent' activation
\end{itemize}
\begin{center}
  \includegraphics[page=44,trim=134 86 134 103,clip,width=254pt]{figure/cs336-l4.pdf}\par
\end{center}
\end{frame}

% CS336 slide 45
\begin{frame}{What else do you need to make DeepSeek MoE v3?}
% hand-typeset (_generator/handwritten/)
\footnotesize
\makebox[\textwidth][c]{\begin{minipage}{400pt}
\textbf{Basic idea:} express the Q, K, V as functions of a lower-dim, `latent' activation
\begin{center}
\includegraphics[page=45,trim=297 211 297 95,clip,width=\linewidth,height=48pt,keepaspectratio]{figure/cs336-l4.pdf}\hspace{14pt}\includegraphics[page=45,trim=439 206 66 100,clip,width=\linewidth,height=48pt,keepaspectratio]{figure/cs336-l4.pdf}
\end{center}
\textbf{Benefits:} when KV-caching, we only need to store $c_t^{KV}$, which can be much smaller.
$W^{UK}$ can be merged into the Q projection\\
\mbox{}\hfill(they also compress queries, for memory savings during training)\quad\raisebox{-0.4\height}{\includegraphics[page=45,trim=561 110 72 239,clip,width=\linewidth,height=20pt,keepaspectratio]{figure/cs336-l4.pdf}}

\textbf{Complexity:} rope conflicts with MLA-style caching.\\[2pt]
\hspace*{1em}Without RoPE -- $\langle Q, K \rangle = \langle hW^Q, W^{UK} c_t^{KV} \rangle = \langle h\,\textcolor{red}{W^Q W^{UK}}, c_t^{KV} \rangle$\\[2pt]
\hspace*{1em}With RoPE -- $\langle QR_q, R_k K \rangle = \langle hW^Q R_q, R_k W^{UK} c_t^{KV} \rangle = \langle h\,\textcolor{red}{W^Q R_q R_k W^{UK}}, c_t^{KV} \rangle$\par\smallskip
{\centering\color{gray} The solution -- Have a few non-latent key dimensions that can be rotated\par}
\end{minipage}}
\end{frame}

% CS336 slide 46
\begin{frame}{What else do you need to make DeepSeek MoE v3?}
% hand-typeset (_generator/handwritten/)
\textbf{MTP:} Have small, lightweight models that predict multiple steps ahead
\begin{columns}[T]
\begin{column}{0.55\textwidth}
\centering
\includegraphics[page=46,trim=39 124 319 114,clip,width=\linewidth]{figure/cs336-l4.pdf}\\
{\scriptsize (But they only do MTP with one token ahead)\hfill [Deepseek v3]\par}
\end{column}
\begin{column}{0.43\textwidth}
\centering
\includegraphics[page=46,trim=413 106 35 97,clip,width=\linewidth,height=72pt,keepaspectratio]{figure/cs336-l4.pdf}\\
{\raggedleft\scriptsize [EAGLE]\par}
\end{column}
\end{columns}
\begin{center}
\includegraphics[page=46,trim=224 4 213 323,clip,width=\linewidth,height=32pt,keepaspectratio]{figure/cs336-l4.pdf}
\end{center}
\vspace{-6pt}
{\raggedleft\footnotesize (See paper for ablations)\par}
\end{frame}

% CS336 slide 47
\begin{frame}{MoE summary}
\begin{itemize}
  \item MoEs take advantage of sparsity -- not all inputs need the full model
  \begin{itemize}
    \item Discrete routing is hard, but top-k heuristics seem to work
  \end{itemize}
  \item Lots of empirical evidence now that MoEs work, and are cost-effective
\end{itemize}
\end{frame}

% FRAMES-END

%%%%%%%%%%%%%%%%%%%%%%%%%%%%%%%%%%%%%%%%%%%%%%%%%%%%%%%%%%%%%%%%%%
\section{Where This Is Going}
%%%%%%%%%%%%%%%%%%%%%%%%%%%%%%%%%%%%%%%%%%%%%%%%%%%%%%%%%%%%%%%%%%

% [O] original bridge frame (draft for the user to rewrite); listed in order.txt
\begin{frame}{Where This Is Going}
\begin{itemize}
  \item The changes in this session act on the Transformer \emph{block}, not on the type of data.
  \item Vision Transformers use them too, e.g.:
  \begin{itemize}
    \item ViT-22B: parallel layers, QK normalization, no biases on the QKV projections and LayerNorms \citebutton{Dehghani et al., 2023}{https://arxiv.org/abs/2302.05442};
    \item V-MoE: a sparse mixture-of-experts Vision Transformer \citebutton{Riquelme et al., 2021}{https://arxiv.org/abs/2106.05974}.
  \end{itemize}
  \item Session 10 takes the Transformer to images, multimodal data and video.
\end{itemize}
\end{frame}


%%%%%%%%%%%%%%%%%%%%%%%%%%%%%%%%%%%%%%%%%%%%%%%%%%%%%%%%%%%%%%%%%%
\section{License \& Attribution}
%%%%%%%%%%%%%%%%%%%%%%%%%%%%%%%%%%%%%%%%%%%%%%%%%%%%%%%%%%%%%%%%%%

\begin{frame}{License \& Attribution}
\begin{itemize}
\item \textbf{Source.} All frames except the four listed below are adapted from \textbf{Stanford CS336 \emph{Language Modeling from Scratch}}, Spring 2025, Lectures 3 (\emph{Architectures, hyperparameters}) and 4 (\emph{Mixture of experts}) by Tatsunori Hashimoto, \url{https://github.com/stanford-cs336/spring2025-lectures} --- \textbf{Copyright 2025 Stanford University}, licensed under the \textbf{MIT License} (full notice in \texttt{\_sources/}). One frame per source slide, in the source's order; text and formulas retypeset in the source's wording, figures cropped from the source PDF.
\item \textbf{Figures} inside the slides are from the papers credited on each frame, reused with that attribution.
\item \textbf{Original to this course.} \enquote{Since 2017: What Changed Inside the Transformer?}, \enquote{Compute vs.\ Memory}, \enquote{The KV Cache}, \enquote{Where This Is Going}.
\end{itemize}
\end{frame}

\endlecture
\end{document}
